% ------------------------------------------------------------------
%  Introduction generale (environ 10 lignes)
% ------------------------------------------------------------------
\section*{Introduction générale}
\phantomsection
\addcontentsline{toc}{section}{Introduction générale}
\markboth{Introduction générale}{}

Le TP Micro-332 a pour but de fabriquer puis de caractériser une résistance électrique en couche mince d'aluminium, en parcourant les étapes de base de la microfabrication \cite{intro}. La résistance d'une piste de section rectangulaire vaut
\begin{equation}
  R = \rho\,\frac{L}{W\,T} = R_\square\,\frac{L}{W},
  \qquad R_\square = \frac{\rho}{T},
  \label{eq:resistance}
\end{equation}
où la longueur $L$ et la largeur $W$ sont fixées par le masque, tandis que la résistivité $\rho$ et l'épaisseur $T$ du film dépendent de la technologie ; $R_\square$ est la résistance par carré.

Deux voies permettent de structurer le métal. En \emph{gravure chimique}, on dépose l'aluminium, on le couvre d'un masque de résine positive, puis on grave les zones non protégées. En \emph{lift-off}, on structure d'abord une résine à profil rentrant (undercut), on dépose le métal par-dessus, puis on dissout la résine, qui emporte le métal situé au-dessus d'elle. Les six sujets couvrent ces deux voies et leur caractérisation (tableau ci-dessous) ; chacun suit la structure demandée dans les consignes \cite{consignes}.

\begin{center}
\small
\captionof*{table}{Objectifs des six sujets du TP.}
\begin{tabularx}{\linewidth}{@{}p{4.6cm}X@{}}
  \toprule
  \textbf{Sujet} & \textbf{Ce que l'on cherche à réaliser} \\
  \midrule
  1 -- Résine positive : dose et adhésion & Calibrer la lithographie : épaisseur de l'AZ1512 en fonction de la vitesse de rotation, courbe de contraste (sensibilité et contraste $\gamma$), effet du HMDS sur l'adhésion \\
  2 -- Lift-off simple couche et bicouche & Obtenir un profil rentrant avec une résine négative nLOF 2070 et avec une bicouche LOR 15C / ECI 3027 ; réaliser un lift-off de Ti/Al sur un échantillon fourni (résine nLOF 2020) \\
  3 -- Observations optiques et MEB & Mesurer l'undercut et l'angle des flancs, et comparer la largeur des pistes d'Al à celle du masque \\
  4 -- Gravure de l'aluminium et alignement & Structurer l'Al par gravure humide et aligner un masque sur des motifs existants \\
  5 -- Salle blanche CMi & Suivre un wafer en lithographie automatisée (EVG 150, MLA 150), graver l'Al dans l'ANP et comparer avec le travail manuel en DLL-MED \\
  6 -- Caractérisation électrique & Mesurer la résistance des pistes et en déduire $R_\square$ et $\rho$ \\
  \bottomrule
\end{tabularx}
\end{center}
